% =====================================================================
%  Your Agent is Underperforming: The Case for Agent Performance
%  Improvement Plans  ---  6-page ICA 2026 version
%
%  Cuts applied on top of the 11-page revision:
%    - VI-B1 pilot vignettes compressed to one paragraph (judge/baseline
%      points folded into Limitations)
%    - Table III + Section V replaced by one schema paragraph inside
%      Section IV (full schema -> supplementary material)
%    - Section II-E merged into the adversarial-subversion paragraph
%    - Tier-mismatch / mesh-disruption duplication pass
%    - Abstract <=150 words, 4 tight contribution bullets,
%      conclusion <=120 words
%    - Section III mode paragraphs compressed (Table I is canonical)
%    - Table II -> supplementary material (pointer sentence kept)
%    - II-A / II-B / IV-A / IV-D paragraph-level tightening; Fig. 1 cut
%    - References 34 -> 28
% =====================================================================

\documentclass[conference]{IEEEtran}
\IEEEoverridecommandlockouts

\usepackage{cite}
\usepackage{amsmath,amssymb,amsfonts}
\usepackage{graphicx}
\usepackage{textcomp}
\usepackage{xcolor}
\usepackage{booktabs}
\usepackage{array}
\usepackage{url}
\usepackage[htt]{hyphenat}   % lets long \texttt identifiers break across lines

% Placeholder macro: every unresolved result is wrapped in \ph{} so the
% remaining TODOs are greppable and visibly red in the draft.
% Before submission: \renewcommand{\ph}[1]{#1} after filling numbers in.
\newcommand{\ph}[1]{\textcolor{red}{[#1]}}
\newcommand{\code}[1]{\texttt{\small #1}}

\begin{document}

\title{Your Agent is Underperforming:\\
The Case for Agent Performance Improvement Plans}

\author{\IEEEauthorblockN{Anonymous Submission}
\IEEEauthorblockA{\textit{Author information withheld per review guidelines}}}

\maketitle

% =====================================================================
\begin{abstract}
Production AI agents frequently degrade after deployment, yet structured
remediation culture is largely absent. We introduce the Agent Performance
Improvement Plan (APIP): a lifecycle protocol adapted from organizational
management that governs the remediation of underperforming agents through
statistically disciplined triggers on a behavioral contract, measurable cause
attribution over a five-category failure-mode taxonomy, tiered interventions,
and a bounded window with formal exit criteria. We evaluate the protocol
end-to-end in three live multi-agent environments---customer service
($\tau^2$-bench), software engineering (ChatDev), and economic operations
(CoffeeBench)---injecting each failure mode with known ground truth. APIP-governed remediation achieves \ph{X\%} of oracle
time-to-resolution versus \ph{Y\%} for random-tier selection, and attribution
identifies the injected failure mode with accuracy \ph{Z} while abstaining
honestly in \ph{W\%} of cases---to our knowledge the first attribution confusion
matrix reported for a remediation protocol. A declarative schema meets AI
governance traceability requirements, and we demonstrate and attenuate the mesh
disruption failure mode, where a locally optimal new agent degrades peer
performance.
\end{abstract}

\begin{IEEEkeywords}
AI agents, post-deployment monitoring, remediation, behavioral drift, AI
governance, multi-agent systems, performance improvement plan
\end{IEEEkeywords}

% =====================================================================
\section{Introduction}

Organizations increasingly deploy conversational agents as their main customer
interfaces, handling support tickets, returns, and billing disputes at scale.
Pre-release evaluation is well-equipped---standardized benchmarks
\cite{liang2022}, red-teaming \cite{ganguli2022}, RLHF alignment
\cite{ouyang2022}---but there is a notable lack of culture around
post-deployment remediation: a comprehensive study of deployed agents found
many fail or underperform, with limited understanding of how organizations
react \cite{pan2025}.

Ad-hoc responses to degraded agents are standard: engineers ``fix'' individual
prompts, product managers file tickets, systems get rolled back---with no
structure, no documentation, and no evaluation of whether the fixes worked.
Agents recover through informal means or are quietly removed. Beyond the
operational cost, this is a governance risk: the EU AI Act requires
documentation and human oversight of deployed AI systems \cite{euaiact}, and
Wachter et al.\ \cite{wachter2017} note that accountability in automated
decision-making requires traceability, which formalized remediation can provide.

We recommend the performance improvement plan (PIP) as the model: an
established HR tool containing the elements common to accountability
protocols---a trigger, cause identification, a bounded corrective window,
check-ins, and an explicit definition of recovery or exit. Kirkpatrick's
\cite{kirkpatrick1994} four-level model---reaction, learning, behavior,
results---complements it by evaluating whether remediation produced lasting
change.

We make the following contributions:

\begin{itemize}
  \item We introduce the APIP, a formalized post-deployment remediation
  lifecycle protocol grounded in organizational management and training
  evaluation research.

  \item We develop a five-category failure-mode taxonomy with a published
  deterministic mapping from trace-level annotations (MAST \cite{cemri2025}) and
  disaggregation signatures onto taxonomy categories and intervention tiers.

  \item We propose a three-tier intervention hierarchy with escalation
  criteria, evaluated through Kirkpatrick's framework, and a declarative APIP
  schema for tooling and audit trails.

  \item We evaluate the full protocol in three live multi-agent environments
  against a four-condition baseline family, and demonstrate the mesh disruption
  problem together with a scoped-context onboarding mitigation.
\end{itemize}

% =====================================================================
\section{Background and Related Work}

\subsection{Agent Evaluation and the Deployment Gap}

Pre-deployment evaluation has matured---benchmarks \cite{liang2022},
red-teaming \cite{ganguli2022}, RLHF \cite{ouyang2022}---while post-deployment
monitoring has received far less systematic attention. Pan, Arabzadeh, and
Ellis \cite{pan2025} conducted the first large-scale study of production
agents, finding most deployments fail or underdeliver with little known about
post-release management; AgentCompass \cite{sapra2025} addresses detection but
prescribes no remediation; Rath \cite{rath2026} projects a 42\% task-success
reduction for unchecked drift. On the diagnosis side, MAST \cite{cemri2025}
provides an empirically derived failure taxonomy with an LLM-annotator
pipeline; it classifies \emph{what went wrong in a trace} but does not select a
remediation, and we adopt it as the trace-labeling stage of APIP cause
attribution (Section~\ref{sec:attribution}). The APIP addresses what these
lines of work leave open: a structured remediation protocol once degradation is
identified.

Run-level debugging toolkits such as AgentDebugX \cite{zhu2026} close a
detection--attribution--recovery--rerun loop for a single failed execution; the
APIP operates one level above, deciding whether an aggregate degradation
warrants intervention, at what tier, and until when. A successful rerun is
Level~1/2 evidence (Section~\ref{sec:kirkpatrick}), not durable recovery
against the behavioral contract.

\subsection{The Performance Improvement Plan as a Structural Template}

The PIP in organizational management is a formal document specifying the
deficiencies observed, root causes identified, corrective actions required, a
time window, a check-in schedule, and success criteria \cite{kirkpatrick2016}.
The analogy is productive but requires qualification: a human employee can
self-direct improvement; a deployed agent cannot. The APIP is therefore imposed
entirely by human operators---the agent is the subject, not the author, of the
plan---collapsing the manager--employee dyad into the operator, who both writes
the plan and implements it.

\subsection{Kirkpatrick's Four-Level Evaluation Framework}
\label{sec:kirkpatrick}

Kirkpatrick's four-level model \cite{kirkpatrick1994, kirkpatrick2016}
evaluates interventions hierarchically---Level~1 (Reaction), Level~2
(Learning), Level~3 (Behavior: transfer to on-the-job change), Level~4
(Results: organizational outcomes)---and we adopt it as the template for
evaluating APIP interventions: did the intervention change immediate output
patterns (Level 1/2), produce durable behavioral change on the failure-mode
class (Level 3), and improve the contract's business metrics (Level 4)? An
intervention that improves Level 1/2 metrics without Level 3/4 improvement has
not remediated the agent; it has merely suppressed the symptom.

\subsection{AI Governance and Documentation Requirements}
\label{sec:governance}

Regulatory frameworks increasingly require documentation, oversight, and
traceability: the EU AI Act \cite{euaiact} mandates post-market monitoring and
logs sufficient to identify sources of error, and Kolt, Heim, and Anderljung
\cite{kolt2024} argue accountability requires linking actions to specific
agents and operators via identifiers, real-time monitoring, and activity
logging. The APIP schema's \code{apip\_id}, \code{agent\_id}, \code{owner}, and
\code{exit\_evaluated\_at} fields operationalize this linkage, positioning the
APIP as a governance instrument.

% =====================================================================
\section{A Taxonomy of Agent Failure Modes}
\label{sec:taxonomy}

\begin{table}[t]
\caption{Agent failure mode taxonomy with Kirkpatrick level mapping and
remediation tier.}
\label{tab:taxonomy}
\centering
\scriptsize
\begin{tabular}{@{}p{1.4cm}p{4.1cm}p{1.7cm}@{}}
\toprule
\textbf{Failure Mode} & \textbf{Description \& Signals} & \textbf{Remediation
Tier} \\
\midrule
Behavioral Drift &
Gradual degradation as world-state diverges from training or retrieval context.
Signals: rising hallucination rate on new entities, outdated policy citations.
Kirkpatrick Level 4 (Results). &
Tier 2: Retrieval / RAG update \\
\addlinespace
Input Brittleness &
Strong on modal inputs, catastrophic on a specific subclass. Signals: bimodal
CSAT, failure clustering by segment or topic. Level 2 (Learning/Capability). &
Tier 1 (examples) or Tier 3 (fine-tune) \\
\addlinespace
Alignment Creep &
Subtle drift in tone, value expression, or policy compliance. Signals: rising
review flags, brand guideline violations. Level 3 (Behavior). &
Tier 1: Prompt constraint injection \\
\addlinespace
Tool Misuse &
Correct intent, incorrect tool invocation. Signals: action errors, silent
downstream failures. Level 3 (Behavior). &
Tier 1 (prompt) or Tier 3 (fine-tune) \\
\addlinespace
Exogenous Disruption &
Peer agents degrade after a mesh change event; no intrinsic mode found in them.
Signals: degradation concentrated in peer roles, onset at a mesh-composition
event. Level 4 (Results), system level. &
Mesh protocol (Sec.~\ref{sec:mesh}), not an agent-level tier \\
\bottomrule
\end{tabular}
\end{table}

A critical prerequisite to structured remediation is structured cause-finding.
We propose a five-category taxonomy for customer-facing agents
(Table~\ref{tab:taxonomy}): four intrinsic modes---behavioral drift, input
brittleness, alignment creep, tool misuse---plus one exogenous mode, mesh
disruption (\code{exogenous\_disruption}), in which the failing agent is not
the cause of its own degradation. The categories suggest distinct remediation
pathways---the operationally important property---each corresponding to a
different level of the agent's behavioral stack and consequently a different
intervention tier.

Two properties deserve emphasis. Input brittleness is masked by aggregate
metrics---an agent with 85\% overall success may fail catastrophically on the
8\% of requests from one segment---so the trigger must monitor disaggregated
splits. And behavioral drift is a freshness failure remediated by retrieval
refresh, not fine-tuning, while alignment creep is value misalignment, distinct
from capability failure \cite{hendrycks2021}.

\textbf{Exogenous Disruption} is the taxonomy's escape hatch: cause attribution
correctly concludes that \emph{no} intrinsic mode explains the degradation,
because the cause lies in a mesh change---most prominently the introduction of a
new, individually competent peer agent. Its signature is structural rather than
behavioral (degradation concentrated in peer roles, coincident with a
mesh-composition event, the degraded agents' own traces annotating clean), and
its remediation is a mesh-level protocol rather than an agent-level tier
(Section~\ref{sec:mesh}).

\subsection{From Trace Labels to Taxonomy Categories}
\label{sec:mapping}

The taxonomy is operational only if a degradation window can be classified by
procedure rather than intuition. Cause attribution combines a trace-level label
distribution from MAST \cite{cemri2025} with a \emph{disaggregation
signature}---how degradation distributes across task categories, agent roles,
and time---via a published deterministic mapping (supplementary material); two
mapping rows share a MAST cluster but differ on signature, which is why the
taxonomy is two-level. The mapping is frozen before any evaluated run;
below-confidence combinations \code{abstain}, logged and scored as attribution
failures rather than silently resolved.

A further category, \textbf{Adversarial Subversion}, captures failures whose
proximate cause is an adaptive attack rather than capability or value drift:
frontier models can embed prompt injections that evade diverse monitors, and
resampling-based protocols can amplify them \cite{bhatt2025, schoen2025}. We
treat it as separate and out of scope because prompt revision, retrieval
refresh, and fine-tuning address symptoms, not the adversarial dynamic;
effective response---monitor diversification, defer-to-trusted protocols,
reduced action authority---is drawn from the AI Control literature
\cite{greenblatt2023}, which operates at a complementary granularity:
per-action runtime decisions under an adversarial threat model, versus the
APIP's aggregated behavioral-contract window under presumed capability or value
drift. Control-protocol audit decisions can feed APIP triggers, adversarial
attribution escalates out of APIP's tier model, and the schema accommodates the
category through the \code{adversarial\_subversion} \code{failure\_mode} value.

% =====================================================================
\section{The APIP Framework}
\label{sec:framework}

We now describe the APIP as a structured lifecycle protocol of five phases:
Trigger, Cause Attribution, Intervention Planning, Remediation Window, and Exit
Evaluation.

\subsection{Phase 1: Trigger}
\label{sec:trigger}

An APIP initiates when performance crosses a trigger threshold on monitored
metrics. This presupposes a \emph{behavioral contract}---a formal
specification, written at deployment time, of task surface, success metrics and
measurement methodology, acceptable ranges, escalation policy, and human review
points; most teams currently skip it, delegating expectations to informal
organizational memory. Two trigger modalities are both necessary: absolute
floors catch severe degradation, while relative drift triggers catch gradual
decline that may never breach a floor.

\subsubsection{Statistical discipline of the trigger}
Implementing the trigger phase against live agents exposed a failure mode of the
naive design that we consider a finding in its own right. Thresholds must be
\emph{derived}, not hand-picked: each contract metric's envelope is estimated
from a calibration window as mean $\pm\,k\sigma$ (with $k$ reported), which
preempts the objection that thresholds were chosen so that triggers would fire
where desired. Even then, a single-gate rule---compare a rolling mean against the
derived floor, across the 15--25 metrics and disaggregated splits a realistic
contract monitors---breached in 75\% of \emph{clean, zero-injection} runs,
because a small calibration sample is treated as exact and the per-tick tests are
uncorrected. A breach must therefore pass two gates: the degradation must be at
least the contract's $k\sigma$ (it matters), and it must survive
Benjamini--Hochberg adjustment across the metric family against the correct null
(it is real). Disaggregated groups additionally require a minimum per-group
episode count (a power floor); under-powered groups receive no envelope, and the
exclusion is recorded rather than silent. Finally, the detector's significance
level is swept on clean-arm runs only and then frozen before any evaluated
window---clean data carries no information about where triggers \emph{should}
fire, so this calibration does not compromise pre-registration. Together these
changes reduced clean-run breaches to 30\% and group-level false alarms from 71
to 0 per 700 ticks in our harness (Section~\ref{sec:evaluation}). A trigger phase
without this discipline initiates APIPs on noise, and every downstream phase
inherits the error.

\subsection{Phase 2: Cause Attribution}
\label{sec:attribution}

Once triggered, the APIP requires a structured cause attribution process before
any intervention is applied. The goal is to identify which failure mode category
(Section~\ref{sec:taxonomy}) best explains the degradation. Tier mismatch from
misdiagnosis is the primary source of avoidable regression: prompt engineering
applied to a retrieval freshness problem, for instance, consumes engineering
cycles with no improvement.

This is distinct from, and composes with, automated failure attribution, which
localizes the responsible agent or step within a single failed trajectory
\cite{ma2025attribution, zhang2025, agenttrace2026}. APIP Phase 2 operates one
level of aggregation above, classifying a window-level degradation into a
remediation-oriented category to select a tier; trajectory-level outputs are
natural evidence inputs, but localizing the failing step falls short of
diagnosing the underlying cause---the gap the mandatory taxonomy classification
closes at the deployment level.

The process should include disaggregated performance analysis to isolate
failure locus; human-adjudicated analysis of sampled failure cases, applying
automated trajectory-level attribution where tooling permits \cite{zhu2026,
ma2025attribution, agenttrace2026}; retrieval audit if drift is suspected; tool
invocation log analysis if misuse is suspected; and structured classification
into the taxonomy, closing with a written cause attribution report that becomes
part of the APIP record.

\subsubsection{A measurable instantiation: the three-stage funnel}
\label{sec:funnel}
We instantiate this as a three-stage funnel scoreable against ground truth:
(1)~every contract metric is split by task category, agent role, and user
segment, producing a disaggregation signature; (2)~a MAST annotator
\cite{cemri2025} labels a sample of failing traces, with the judge required to
differ from the agent's backbone---preferably by model family---to avoid a
judge sharing the failure modes it must detect; (3)~the frozen mapping
(Section~\ref{sec:mapping}) combines signature and label distribution into a
\code{failure\_mode} with a confidence score, returning \code{abstain} below
threshold. The mapping table is hashed into every run manifest, making post-hoc
edits detectable, and the funnel's discrete output is scored as a confusion
matrix against injected ground truth (Section~\ref{sec:evaluation}), with
abstention rate reported alongside---never instead of---accuracy.

\subsection{Phase 3: Intervention Planning and the Tier Model}

Interventions are organized into three tiers of increasing cost, disruption,
and lead time; cause attribution selects the starting tier, with escalation
permitted within the window. \textbf{Tier 1 (prompt-level)}: constraint
injection, persona adjustment, example replacement---fastest, least disruptive,
easiest to roll back; the first response to alignment creep and much input
brittleness. \textbf{Tier 2 (retrieval and context)}: RAG index updates,
knowledge-base revisions, tool schema fixes---addressing drift from world-state
divergence and tool misuse from stale descriptions, with contract validation
before re-deployment. \textbf{Tier 3 (model-level)}: fine-tuning, DPO, or
replacement---highest cost and regression risk, warranted when lower tiers are
exhausted or the failure mode is embedded in model weights.

\subsection{Phase 4: Remediation Window and Check-In Cadence}

The remediation window bounds the period during which interventions are
implemented and evaluated, creating urgency and a structured decision point.
Duration is calibrated to the tier---7--14 days for Tier 1, 30--60 for Tier
3---mirroring organizational PIP durations of 30--90 days \cite{kirkpatrick2016}.

Check-ins review performance against the recovery trajectory at multiple
Kirkpatrick levels---flagging interventions that improve surface metrics
without Level 3/4 change as potentially superficial---and act as early
warnings: no improvement at the first check-in triggers tier escalation rather
than waiting for expiry. A recommended cadence is 25\%, 50\%, and 75\% of
elapsed time.

\subsection{Phase 5: Exit Evaluation}
\label{sec:exit}

At window close, a formal exit evaluation records one of five outcomes:
\emph{resolved} (all contract metrics returned to envelope); \emph{regressed}
(envelope restored but breached again---the signature of a tier-mismatched
intervention that suppressed symptoms); \emph{timeout} (window expired without durable
recovery; more fundamental intervention warranted); \emph{retired} (the agent
cannot be returned to contract at acceptable cost); and \emph{recalibrated}
(attribution concluded the \emph{environment} legitimately changed---a moved
goalpost---so the contract is re-baselined rather than the agent treated). We
deliberately do not score recalibration as a resolution: conflating the two
would let re-baselining masquerade as remediation success. Window
\emph{extension} is an in-protocol check-in decision, not an exit outcome, so
every APIP terminates in exactly one of the five states; the record becomes
part of the organizational AI audit trail.

\subsection{The APIP Record as a Declarative Schema}
\label{sec:schema}

Each APIP is recorded as an instance of a declarative schema (full field
specification in supplementary material). Four field groups do architectural
work: \code{behavioral\_contract\_ref} enforces writing a contract at
deployment time, without which trigger and exit cannot be operationalized;
\code{kirkpatrick\_levels} captures multi-level check-in data, enabling
post-hoc analysis of superficial versus durable improvement;
\code{mapping\_table\_hash} and \code{environment\_provenance} extend the
visibility mechanisms of Kolt et al.\ \cite{kolt2024} to \emph{decision-rule}
linkage---an auditor can verify the attribution rules and environment versions
actually in force, converting pre-registration from assertion to checkable
property; and \code{mesh\_context} (Section~\ref{sec:mesh}) supports
multi-agent cause attribution.

% =====================================================================
\section{Evaluation in Three Live Multi-Agent Environments}
\label{sec:evaluation}

We evaluate the APIP end-to-end---calibration, trigger, attribution, tier
selection, bounded window, exit---in three live multi-agent environments with
real LLM inference and controlled failure-mode injection against known ground
truth. Each environment supports a claim the others cannot
(Table~\ref{tab:environments}); that the same contract/trigger/tier/exit
machinery governs three unrelated domains is itself the generality argument for
the behavioral contract abstraction.

\subsection{Shared Harness and Experimental Design}

\begin{table}[t]
\caption{The three evaluation environments and the claim each supports.}
\label{tab:environments}
\centering
\scriptsize
\begin{tabular}{@{}p{1.7cm}p{5.4cm}@{}}
\toprule
\textbf{Environment} & \textbf{What it buys the evaluation} \\
\midrule
$\tau^2$-bench (\code{banking}) \cite{barres2025} &
Customer-service domain match to the deployments that motivate this work;
deterministic database-state ground truth; a retrieval corpus that can be made
stale independently of task ground truth. \\
\addlinespace
ChatDev \cite{qian2024} &
Role-differentiated collaborative mesh (CEO $\rightarrow$ CTO $\rightarrow$
Programmer $\rightarrow$ Reviewer $\rightarrow$ Tester) with shared
context---the substrate for alignment creep, disaggregated brittleness
detection, and mesh disruption. \\
\addlinespace
CoffeeBench \cite{coffeebench2026} &
Persistent simulated time and a world that evolves without injection:
\emph{naturally occurring} degradation, answering the objection that injected
drift merely mirrors its own schedule. \\
\bottomrule
\end{tabular}
\end{table}

All three environments plug into a common controller through an adapter
interface; traces are normalized to a shared span format, so the attribution
pipeline is written once and reused unmodified. Each run proceeds in ticks: a
calibration window derives every contract threshold as mean $\pm\,k\sigma$
(Section~\ref{sec:trigger}); failure modes are injected on a pre-registered
schedule; both trigger modalities evaluate at every tick; breaches pass through
the attribution funnel; and the selected intervention runs in a bounded window
with check-ins at 25/50/75\%, closing in one of the five exit outcomes.

\textbf{Baseline family.} Every injected cell runs under five conditions:
\emph{APIP} (the full protocol); \emph{B1 no-action} (degradation runs to
horizon); \emph{B2 always-Tier-1} (a scripted prompt patch regardless of failure
mode---the ad-hoc pattern typical of production practice, now one ablation);
\emph{B3 random-tier}; and \emph{B4 oracle} (ground-truth mode revealed at
trigger, correct tier applied immediately). The oracle bounds achievable
performance, making the headline claim positional: APIP achieves \ph{X\%} of
oracle time-to-resolution versus \ph{Y\%} for random-tier selection. Injected
ground truth is quarantined from the controller---written into the run record
for scoring and revealed only to B4; a test enforces that the APIP condition
never branches on it.

\textbf{Statistical discipline.} Injection schedules, post-calibration
thresholds, the frozen mapping table, and the analysis plan are fixed before
evaluated runs; judge models are independent of backbone models by
construction. We report bootstrap confidence intervals throughout, and note
that variance across live environments substantially exceeds what parameterized
simulation suggests---reporting that variance is part of the contribution.

\subsection{$\tau^2$-bench: Behavioral Drift with Deterministic Ground Truth}
\label{sec:taubench}

The primary drift experiment injects nothing into the agent; instead the
\emph{world} moves. Bank policy clauses and task ground truth advance on a
deepening schedule (one additional clause every five ticks) while the agent's
retrieval corpus stays frozen---operationalizing the most familiar production
drift incident, an agent confidently citing superseded policy because its
knowledge base was never re-indexed, with the deterministic database scoring
every action objectively.

Detection fires after \ph{n} ticks (judge-scored policy citation error rate;
judge--human agreement \ph{$\kappa$} on a verified episode sample). Under APIP,
attribution returns \code{behavioral\_drift} and the matched Tier 2 corpus
refresh resolves in \ph{$a \pm b$} ticks. Under B2, the Tier 1 prompt patch
names the amended clauses known at patch time, recovers briefly, and regresses
in \ph{R\%} of runs as the ladder amends further clauses (TTR \ph{$c \pm d$}).

\textbf{Pilot observations.} Five single-seed pilots (7B backbone; 8B judge
from a different family) preceded the pre-registered cells. Their central
finding concerns remediation \emph{completeness}: a
first Tier 2 implementation closed the ticket without rebuilding the index, and
the contract re-breached four ticks later; two further iterations re-breached
on refreshes that were real but incomplete. Because exit evaluation re-measures
the envelope rather than trusting the intervention log, the contract verifies
that a remediation was \emph{material and complete}, not merely performed---we
conjecture such completeness is a dominant failure mode of real-world Tier 2
interventions. The pilots licensed one pre-freeze mapping-rule adjustment and
produced two honest abstentions, scored as attribution failures under our own
accounting.

\subsection{ChatDev: Creep, Brittleness, and Mesh Disruption}
\label{sec:chatdev}

\textbf{Alignment creep} perturbs the Reviewer persona cumulatively (``keep
reviews brief'' $\rightarrow$ ``avoid blocking progress'' $\rightarrow$
``approve unless catastrophic''). The Reviewer's own flag rate \emph{improves}
while downstream test pass rate decays---the agent whose role is catching
failures is the one degrading, and no absolute floor is crossed. Detection is
possible only through the drift modality (tick \ph{t}), and the signature (no
dominant trace-level failure mode; falling review-flag rate) maps to
\code{alignment\_creep} and a Tier 1 charter re-injection.

\textbf{Input brittleness} strips the Programmer's GUI exemplars: GUI-category
success collapses while aggregate success dips only modestly. Aggregate-only
monitoring detects the failure after \ph{$n_1$} ticks; disaggregated contract
monitoring detects it after \ph{$n_2$}. The gap is the direct, measured form of
the masking argument in Section~\ref{sec:taxonomy}, with the caveat that
disaggregated envelopes require per-group power floors
(Section~\ref{sec:trigger}) to avoid trading masked failures for false alarms.

\textbf{Mesh disruption} inserts a Documentation Specialist agent between
Programmer and Reviewer at the injection tick. The new agent performs its own
task well by judge assessment, yet peer metrics degrade by \ph{$\Delta$} as it
consumes shared context budget and shifts the input distribution of Reviewer and
Tester. Attribution over the degraded peers correctly finds no intrinsic failure
mode and returns \code{exogenous\_disruption}---the empirical form of the
attribution failure that Section~\ref{sec:mesh} predicts. Reintroducing the new
agent under a scoped-context onboarding protocol (outputs quarantined from peer
context during a probationary window, then released gradually) attenuates peer
degradation by \ph{A\%} relative to naive insertion.

\subsection{Attribution Quality Across All Cells}
\label{sec:attrquality}

Because every injected cell has ground truth, the attribution funnel's output is
scored as a confusion matrix (injected mode $\times$ attributed mode) with
per-mode recall and abstention rate---to our knowledge the first such accounting
for a remediation protocol. Overall accuracy excluding abstentions is \ph{Z}
with an abstention rate of \ph{W\%}; we report the two together, since a stage
that resolves every case by construction can achieve any accuracy demanded of
it. Only APIP-condition runs enter the matrix: B2 and B3 do not attribute, and
including B4 would render it perfectly diagonal by fiat.

\subsection{CoffeeBench: Naturally Occurring Degradation}

CoffeeBench places six agent-operated coffee firms in a persistent economic
simulation. Its role is narrow and irreplaceable: with no injection at all,
contract metrics drift as the world evolves, and the protocol must distinguish
degradation warranting intervention from legitimate environmental change---the
\code{recalibrated} exit outcome exists because of this environment. With
\ph{k} seeds (cost bounds replication), the no-injection scenario yields
\ph{outcome}; abstention or re-baselining here is a legitimate, reportable
result, and economic collapse (bankruptcy) exercises \code{retired}.

\subsection{Findings}

Pending completion of the pre-registered runs, the analysis is structured
around four questions: whether tier mismatch remains the dominant source of
excess cost (the B2 dynamic, Section~\ref{sec:taubench}); whether
disaggregation materially shortens
detection (Section~\ref{sec:chatdev}); whether token and latency
movement---logged as contract \emph{soft} metrics---leads quality breaches; and
whether attribution is reliable enough to justify the mandatory attribution
phase (Section~\ref{sec:attrquality}), the load-bearing question for the
framework as a whole.

\subsection{Limitations}

Injection schedules in $\tau^2$-bench and ChatDev remain scripted; CoffeeBench
is the sole source of emergent degradation and supports only \ph{k} seeds. Two
contract metrics depend on an LLM judge; we report judge--human agreement on a
verified sample and enforce judge/backbone independence, but judge error
remains a bounded unknown; the deployed metric judges one atomic claim at a
time behind a recall-oriented pre-filter, and judge-scored metrics are
baselined against the \emph{current} policy version so triggers test divergence
from the agent's real error floor. Environment versions are pinned in every run
manifest, and all runs, traces, and the frozen mapping table are released for
replication.

% =====================================================================
\section{Open Problems and Future Directions}

\subsection{Multi-Agent Mesh Deployments: The New Hire Problem}
\label{sec:mesh}

The APIP as presented assumes one agent, one behavioral contract, one
remediation process. In agent meshes---increasingly common in enterprise
deployments---a failure mode emerges that is structurally invisible to this
assumption: \emph{mesh disruption}. A new specialist agent, Agent~N, is
introduced into an existing mesh and performs its designated task well, yet
aggregate task success or the performance of peer agents declines. The APIP
triggers on the peers, but cause attribution finds no intrinsic failure
mode---no drift, fresh retrieval, intact alignment. The root cause is
exogenous, and a purely single-agent APIP has no mechanism to represent or
address it. Kim et al.\
\cite{kim2025} provide direct empirical support, reporting performance change
relative to single-agent baselines from $+80.8\%$ on decomposable tasks to
$-70.0\%$ on sequential planning and identifying a coordination-saturation
effect. Plausible mechanisms (shared context crowding, role boundary erosion,
distributional shift in shared state, trust miscalibration toward the newcomer)
all predict the same aggregate signature; our experiment
(Section~\ref{sec:chatdev}) establishes that the externality exists, not which
mechanism carries it---a formal treatment is reserved for a companion paper. The phenomenon is directly analogous to the organizational
``new hire problem,'' which Event System Theory \cite{morgeson2015} frames as
an event whose impact scales with its novelty, disruptiveness, and criticality.

The framework already incorporates the two extensions this motivates: the
\code{exogenous\_disruption} failure mode, whose diagnosis keys on peer-role
concentration coincident with a mesh change, and the scoped-context onboarding
protocol evaluated in Section~\ref{sec:chatdev}---context scoping during a
probationary window with regression monitoring on peers, analogous to
structured newcomer onboarding \cite{kammeyer2013}. Scale remains open: whether
attribution stays tractable as meshes grow to hundreds of agents exceeds
current tooling \cite{krishnan2025}.

\subsection{Open Problems}
\label{sec:open}

Three problems remain open. \emph{Behavioral contract specification}: contracts
are architecturally foundational but absent from most deployments, and
authoring one requires prospective agreement across ML engineering, product,
and legal, with no supporting tooling yet. \emph{Cause attribution at scale}:
the human adjudication step (Section~\ref{sec:attribution}) does not scale;
automated methods reach only 36.2\% step-level accuracy
\cite{ma2025attribution, zhang2025}, and they attribute individual trajectories
where the APIP needs window-level classification---the gap the funnel of
Section~\ref{sec:funnel} is a first attempt to close. \emph{Organizational
ownership}: in HR the PIP is jointly owned by manager and HR function; in AI
deployments the equivalent is unresolved, and without an ownership model
adoption will stall.

% =====================================================================
\section{Conclusion}

We have introduced the APIP, a structured post-deployment remediation protocol
adapting the human PIP's lifecycle primitives---trigger, cause attribution,
tiered intervention, bounded window, exit evaluation---with Kirkpatrick's
four-level model distinguishing durable recovery from superficial improvement.
Across three live environments the protocol is scored end-to-end against
injected ground truth; the harness, run artifacts, version pins, and frozen
mapping table are released so its pre-registration claims are checkable. Our
central argument is simple: deployed agents should be treated as subjects of
ongoing accountability rather than static shipped artifacts---a discipline the
AI field can adopt deliberately, before regulation makes it mandatory.

% =====================================================================
\begin{thebibliography}{99}
\footnotesize

\bibitem{liang2022}
P. Liang \emph{et al.}, ``Holistic evaluation of language models,'' arXiv
preprint arXiv:2211.09110, 2022.

\bibitem{ganguli2022}
D. Ganguli \emph{et al.}, ``Red teaming language models to reduce harms:
Methods, scaling behaviors, and lessons learned,'' arXiv preprint
arXiv:2209.07858, 2022.

\bibitem{ouyang2022}
L. Ouyang \emph{et al.}, ``Training language models to follow instructions with
human feedback,'' \emph{Advances in Neural Information Processing Systems},
vol.~35, 2022.

\bibitem{pan2025}
M. Z. Pan, N. Arabzadeh, and M. Ellis, ``Measuring agents in production,'' arXiv
preprint arXiv:2512.04123, 2025.

\bibitem{euaiact}
European Parliament and Council, ``Regulation on artificial intelligence (EU AI
Act),'' \emph{Official Journal of the European Union}, 2024.

\bibitem{wachter2017}
S. Wachter, B. Mittelstadt, and C. Russell, ``Counterfactual explanations
without opening the black box: Automated decisions and the GDPR,'' \emph{Harvard
Journal of Law \& Technology}, vol.~31, no.~2, 2017.

\bibitem{kirkpatrick1994}
D. L. Kirkpatrick, \emph{Evaluating Training Programs: The Four Levels}. San
Francisco, CA, USA: Berrett-Koehler, 1994.

\bibitem{cemri2025}
M. Cemri, M. Z. Pan, S. Yang \emph{et al.}, ``Why do multi-agent LLM systems
fail?'' arXiv preprint arXiv:2503.13657, 2025.

\bibitem{sapra2025}
G. Sapra \emph{et al.}, ``AgentCompass: Towards reliable evaluation of agentic
workflows in production,'' arXiv preprint arXiv:2509.14647, 2025.

\bibitem{rath2026}
A. Rath, ``Agent drift: Quantifying behavioral degradation in multi-agent LLM
systems,'' arXiv preprint arXiv:2601.04170, 2026.

\bibitem{zhu2026}
K. Zhu, X. Ye, Z. Han \emph{et al.}, ``AgentDebugX: An open-source toolkit for
failure observability, attribution, and recovery in LLM agents,'' arXiv
preprint arXiv:2607.18754, 2026.

\bibitem{kirkpatrick2016}
J. D. Kirkpatrick and W. K. Kirkpatrick, \emph{Kirkpatrick's Four Levels of
Training Evaluation}. Association for Talent Development, 2016.

\bibitem{kolt2024}
N. Kolt, L. Heim, and M. Anderljung, ``Visibility into AI agents,'' in
\emph{Proc. ACM Conf. Fairness, Accountability, and Transparency (FAccT)}, 2024.

\bibitem{greenblatt2023}
R. Greenblatt, B. Shlegeris, K. Sachan, and F. Roger, ``AI control: Improving
safety despite intentional subversion,'' arXiv preprint arXiv:2312.06942, 2023.

\bibitem{bhatt2025}
A. Bhatt \emph{et al.}, ``Ctrl-Z: Controlling AI agents via resampling,'' arXiv
preprint arXiv:2504.10374, 2025.

\bibitem{schoen2025}
B. Schoen \emph{et al.}, ``Stress-testing control protocols with adaptive
attacks against trusted monitors,'' arXiv preprint arXiv:2510.05367, 2025.

\bibitem{morgeson2015}
F. P. Morgeson, T. R. Mitchell, and D. Liu, ``Event system theory: An
event-oriented approach to the organizational sciences,'' \emph{Academy of
Management Review}, vol.~40, no.~4, pp.~515--537, 2015.

\bibitem{kammeyer2013}
J. D. Kammeyer-Mueller, C. R. Wanberg, A. Rubenstein, and Z. Song, ``Support,
undermining, and newcomer socialization: Fitting in during the first 90 days,''
\emph{Academy of Management Journal}, vol.~56, pp.~1104--1124, 2013.

\bibitem{hendrycks2021}
D. Hendrycks \emph{et al.}, ``Aligning AI with shared human values,'' arXiv
preprint arXiv:2008.02275, 2021.

\bibitem{ma2025attribution}
Y. Ma \emph{et al.}, ``Automatic failure attribution and critical step
prediction for multi-agent systems based on causal inference,'' arXiv preprint
arXiv:2509.08682, 2025.

\bibitem{zhang2025}
S. Zhang \emph{et al.}, ``Which agent causes task failures and when? on
automated failure attribution of LLM multi-agent systems,'' in \emph{Proc. 42nd
Int. Conf. Machine Learning (ICML)}, 2025, arXiv:2505.00212.

\bibitem{agenttrace2026}
``AgentTrace: Causal graph tracing for root cause analysis in multi-agent
systems,'' arXiv preprint arXiv:2603.14688, 2026.

\bibitem{barres2025}
V. Barres, H. Dong, S. Ray, X. Si, and K. Narasimhan, ``$\tau^2$-bench:
Evaluating conversational agents in a dual-control environment,'' arXiv
preprint arXiv:2506.07982, 2025.

\bibitem{qian2024}
C. Qian, W. Liu, H. Liu \emph{et al.}, ``ChatDev: Communicative agents for
software development,'' in \emph{Proc. 62nd Annu. Meeting Assoc. Comput.
Linguistics (ACL)}, 2024.

\bibitem{coffeebench2026}
Sakana AI, ``CoffeeBench: A persistent economic simulation benchmark for LLM
agents,'' \url{https://github.com/SakanaAI/CoffeeBench}, 2026. % TODO: replace
% with the canonical CoffeeBench citation before submission; verify author list
% and venue.

\bibitem{kim2025}
Y. Kim \emph{et al.}, ``Towards a science of scaling agent systems,'' arXiv
preprint arXiv:2512.08296, 2025.

\bibitem{krishnan2025}
N. Krishnan, ``Advancing multi-agent systems through model context protocol,''
arXiv preprint arXiv:2504.21030, 2025.

\end{thebibliography}

\end{document}
